% =====================================================================
%  PLANTILLA: actividad / cuestionario / entrega corta
%  (cabecera compacta en lugar de portada, sin índice)
%  Máster Universitario en Ingeniería Informática (IA) - UPO 2026/2027
% =====================================================================
\documentclass[cabecera,sinindice]{upo-master}

\asignatura{Nombre de la asignatura}
\tipodocumento{Actividad 1}
\title{Título de la actividad}
\author{Javier Gómez López}
% \profesor{Nombre del profesor}
% \rutaLogo{../../Plantilla LaTeX/img/logo-upo}

% Uso de IA: una línea por herramienta (o ninguna si no se ha usado).
\herramientaIA{Claude (Anthropic)}{Claude Opus 5.5}{Revisión de la redacción.}

\begin{document}

\maketitle

\ejercicio[Título del ejercicio]
\begin{enunciado}
Texto de la pregunta tal y como aparece en el enunciado.
\end{enunciado}
\begin{respuesta}
Respuesta razonada.
\end{respuesta}

\ejercicio
\begin{enunciado}
Segunda pregunta.
\end{enunciado}
\begin{respuesta}
Respuesta. Si hay fórmulas: $\E[X] = \sum_x x\,P(X=x)$.
\end{respuesta}

% La declaración de uso de IA es obligatoria: aquí o, si se omite,
% se añade sola al final.
\declaracionIA

% Si hay referencias, descomenta:
% \bibliographystyle{unsrt-es}
% \bibliography{bibliografia}

\end{document}
